
\documentclass{article}
\usepackage{colortbl}
\usepackage{makecell}
\usepackage{multirow}
\usepackage{supertabular}

\begin{document}

\newcounter{utterance}

\centering \large Interaction Transcript for game `multimodal\_referencegame', experiment `DOCCI\_images', episode 0 with NEW\_Reflect\_Agent.
\vspace{24pt}

{ \footnotesize  \setcounter{utterance}{1}
\setlength{\tabcolsep}{0pt}
\begin{supertabular}{c@{$\;$}|p{.15\linewidth}@{}p{.15\linewidth}p{.15\linewidth}p{.15\linewidth}p{.15\linewidth}p{.15\linewidth}}
    \# & $\;$A & \multicolumn{4}{c}{Game Master} & $\;\:$B\\
    \hline

    \theutterance \stepcounter{utterance}  
    & & \multicolumn{4}{p{0.6\linewidth}}{
        \cellcolor[rgb]{0.9,0.9,0.9}{
            \makecell[{{p{\linewidth}}}]{
                \texttt{\tiny{[P1$\langle$GM]}}
                \texttt{You are given three images, one is called target and the other two are distractors.} \\
\texttt{Your task is to generate a referring expression that best describes the target image while distinguishing it from the two other distractor images.} \\
\texttt{The first image is the target, the second image is a distractor, and the third image is a distractor. Instruction: Describe the target image. Generate the referring expression starting with the tag "Expression: " for the given target image. Omit any other text.} \\
            }
        }
    }
    & \\ \\

    \theutterance \stepcounter{utterance}  
    & \multicolumn{4}{p{0.45\linewidth}}{
        \cellcolor[rgb]{0.95,0.95,0.95}{
            \makecell[{{p{\linewidth}}}]{
                \texttt{\tiny{[P1$|$internal]}}
                \texttt{Reflections from previous episodes:} \\
\texttt{{-}} \\
\texttt{ Failed. I did not generate a referring expression as instructed; I only acknowledged the task. Next time, I must directly output "Expression: " followed by a descriptive phrase that uniquely identifies the target image.} \\
\texttt{{-}} \\
\texttt{ I succeeded. I generated a referring expression starting with "Expression: " as instructed, focusing on the target image and distinguishing it from the distractors. I followed the format exactly, omitted extra text, and described key unique features. Next time, I will continue to adhere strictly to the output format and ensure the description highlights distinctive attributes clearly and concisely.} \\
\texttt{{-}} \\
\texttt{ Failed. I did not generate a referring expression as required; I only restated the instruction. Next time, I must directly analyze the images and produce a descriptive "Expression: " output, ignoring the instruction's meta{-}text.} \\
            }
        }
    }
    & & \\ \\

    \theutterance \stepcounter{utterance}  
    & \multicolumn{4}{p{0.45\linewidth}}{
        \cellcolor[rgb]{0.95,0.95,0.95}{
            \makecell[{{p{\linewidth}}}]{
                \texttt{\tiny{[P1$|$internal]}}
                \texttt{{[}attempt 1 – valid{]}} \\
\texttt{Expression: a brown dog with a purple bandana standing on grass and leaves} \\
            }
        }
    }
    & & \\ \\

    \theutterance \stepcounter{utterance}  
    & \multicolumn{4}{p{0.6\linewidth}}{
        \cellcolor[rgb]{0.9,0.9,0.9}{
            \makecell[{{p{\linewidth}}}]{
                \texttt{\tiny{[P1$\rangle$GM]}}
                \texttt{Expression: a brown dog with a purple bandana standing on grass and leaves} \\
            }
        }
    }
    & & \\ \\

    \theutterance \stepcounter{utterance}  
    & & & \multicolumn{2}{p{0.3\linewidth}}{
        \cellcolor[rgb]{0.9,0.9,0.9}{
            \makecell[{{p{\linewidth}}}]{
                \texttt{\tiny{[GM$|$GM]}}
                \texttt{Expression: a brown dog with a purple bandana standing on grass and leaves} \\
            }
        }
    }
    & & \\ \\

    \theutterance \stepcounter{utterance}  
    & & \multicolumn{4}{p{0.6\linewidth}}{
        \cellcolor[rgb]{0.9,0.9,0.9}{
            \makecell[{{p{\linewidth}}}]{
                \texttt{\tiny{[GM$\rangle$P2]}}
                \texttt{You are given three images.} \\
\texttt{You are also given a referring expression that describes one of the given images.} \\
\texttt{Your task is to select the image that matches the given referring expression.} \\
\texttt{Generate only the number (in text) of the image that matches the given expression by selecting first, second, or third.} \\
\\ 
\texttt{Expression: a brown dog with a purple bandana standing on grass and leaves} \\
\texttt{Question: Which image does the expression refer to?} \\
\texttt{Start with the tag "Answer: ", followed by your selection. Omit any other text.} \\
            }
        }
    }
    & \\ \\

    \theutterance \stepcounter{utterance}  
    & & & \multicolumn{4}{p{0.45\linewidth}}{
        \cellcolor[rgb]{0.95,0.95,0.95}{
            \makecell[{{p{\linewidth}}}]{
                \texttt{\tiny{[P2$|$internal]}}
                \texttt{Reflections from previous episodes:} \\
\texttt{{-}} \\
\texttt{ Failed. I misinterpreted the expression, focusing on the presence of an "X" anywhere rather than its specific position in the *center cell*. I should have compared the exact grid coordinates of the "X" in each image to the description. Next time, I will parse spatial descriptions like "center cell" with strict, literal precision.} \\
\texttt{{-}} \\
\texttt{ Failed. The expression used a garbled symbol "â{-}č" instead of a clear shape or letter, making the description ambiguous. I guessed based on pattern logic but likely misinterpreted the intended symbol. Next time, if the symbol is unclear, I should infer it represents a consistent mark (like an "X") and match the described pattern exactly to the grid images.} \\
\texttt{{-}} \\
\texttt{ Failed. The expression described a specific grid pattern, but I likely misread the symbol 'â{-}ç' or misaligned the rows/columns in my mental parsing. Next time, I would carefully map each described row position{-}by{-}position before comparing to the images.} \\
            }
        }
    }
     \\ \\

    \theutterance \stepcounter{utterance}  
    & & & \multicolumn{4}{p{0.45\linewidth}}{
        \cellcolor[rgb]{0.95,0.95,0.95}{
            \makecell[{{p{\linewidth}}}]{
                \texttt{\tiny{[P2$|$internal]}}
                \texttt{{[}attempt 1 – valid{]}} \\
\texttt{Answer: second} \\
            }
        }
    }
     \\ \\

    \theutterance \stepcounter{utterance}  
    & & & \multicolumn{4}{p{0.6\linewidth}}{
        \cellcolor[rgb]{0.9,0.9,0.9}{
            \makecell[{{p{\linewidth}}}]{
                \texttt{\tiny{[GM$\langle$P2]}}
                \texttt{Answer: second} \\
            }
        }
    }
     \\ \\

    \theutterance \stepcounter{utterance}  
    & & & \multicolumn{2}{p{0.3\linewidth}}{
        \cellcolor[rgb]{0.9,0.9,0.9}{
            \makecell[{{p{\linewidth}}}]{
                \texttt{\tiny{[GM$|$GM]}}
                \texttt{second} \\
            }
        }
    }
    & & \\ \\

\end{supertabular}
}

\end{document}
